\section{10H-N reconstruction: multihorizon refinement}
Starting from the 20 reconstructed M libraries, each donor-global plant and
two components was refined with an equal-weight sensor-normalized measured-output
loss at horizons 1, 5, 20 and 50. Component training memberships were held fixed
at their M values, an explicit reconstruction choice. Derivatives included both
the KF recursion and future open-loop forecast. Bounds, covariances and the
eight-coordinate physical coupling were unchanged.

Of 60 proposed refinements, 59 passed success, normalized KKT, feasibility and
training-loss checks; the remaining plant retained its original coefficients.
Validation selected the actual refined bundle for 16 of 20 libraries. At horizon
50, mean paired recipient gains relative to M were 6.49\% for
$Q_{\rm fit}=0.01$ and 26.99\% for $Q_{\rm fit}=1$. Short-horizon gains were mixed:
at horizon 1 they were $-4.04\%$ and $-3.55\%$, and at horizon 5 they were
3.26\% and $-11.16\%$. Individual long-horizon harms remain substantial. These
fresh development results support a feedback diagnostic, not a uniform transfer
claim. The libraries and prefix-only recipient assignments are frozen before
closed-loop evaluation; no controller costs or bounds are retuned.
